\scanpage{102}
\mspage{95}

On vérifie aussi immédiatement que

\[
 \underset{\alpha\in\widetilde C}{\varinjlim} Z_\alpha=\varinjlim X_i=X
\]

donc, $\widetilde C_M$ étant cofinal,

\[
 \underset{\alpha\in\widetilde C_M}{\varinjlim} Z_\alpha=X.
\]

\deleted{Il reste bien entendu que le syst. inductif $Z_\alpha$ sur $\widetilde C$ tout entier, donc \unsure{restreint} sur $\widetilde C_M$,}

Au total, on est arrivé à représenter $X$ sous la forme

\[
 X=\underset{\alpha\in A}{\varinjlim} X_\alpha
\]

\unsure{où|mais} cette fois les $X_\alpha$ sont \uline{dans $M^\pi$} (pas seulement dans $M$), et $A$ est \deleted{\unsure{bien}} grand devant $\pi$. \deleted{adapté à $\pi$ (pas seulement grand devant $\pi$)} Mais on ne \unsure{sait} pas que \deleted{les} toute partie \inserted{filtrante $B$} de cardinal $\leq\pi$ de $A$ admet une borne sup. \inserted{$\alpha_B$} dans $A$, et encore moins que $X_{\alpha_B}=\varinjlim_{\alpha\in B} X_\alpha$.

% slanted
\margin{Si, \unsure{on suppose|en supposant} $(\unsure{a_\pi|a\,\pi})$ c'est ce que \unsure{ns} voulions prouver \unsure{dans} a), que $M$ est stable par $\varinjlim$ grandes devant $\pi$.}

Donc on ne peut \inserted{encore} appliquer directement b). S'il ne s'agissait que de prouver \inserted{la deuxième moitié de} a), \unsure{où on|on en} suppose déjà des $X_i\in M^\pi$ et $I$ grand devant $I$, on n'a rien gagné du tout \unsure{par|pour} nos labeurs !

% archive sheet number: 101
